% Part: lambda-calculus
% Chapter: introduction
% Section: overview

\documentclass[../../../include/open-logic-section]{subfiles}

\begin{document}

\olfileid{lam}{int}{ovr}
\olsection{সারসংক্ষেপ}

অ্যালোনজো চার্চ ১৯৩০-এর দশকের গোড়ায় প্রথমে গঠনমূলক যুক্তিবিদ্যার ভিত্তি
হিসেবে ল্যাম্বডা ক্যালকুলাস নির্মাণ করেছিলেন, গণনসাধ্য অপেক্ষকগুলির মডেল
হিসেবে \emph{নয়}। কিন্তু অচিরেই দেখানো হয় যে এটি গণনযোগ্যতার অন্যান্য
সংজ্ঞার সমতুল্য, যেমন টুরিং-গণনসাধ্য অপেক্ষক এবং আংশিক পুনরাবৃত্ত অপেক্ষক।
প্রথমে এটি কিছুটা বিস্ময়কর বলে মনে হয়েছিল—এই তথ্যটিই বৈশিষ্ট্যায়নটিকে আরও
আকর্ষণীয় করে তোলে।

কোনো অপেক্ষকের নাম সংজ্ঞায়িত না করে, তাকে সংজ্ঞায়িত করে এমন একটি সাংকেতিক
রাশির সাহায্যে সরাসরি সেই অপেক্ষককে নির্দেশ করার সুবিধাজনক উপায় হলো
ল্যাম্বডা-লিপি। ``$f$-কে $f(x) = x + 3$ দ্বারা সংজ্ঞায়িত অপেক্ষক ধরা যাক''
না বলে বলা যায়, ``$f$-কে $\lambd[x][(x + 3)]$ অপেক্ষক ধরা যাক।'' অন্যভাবে
বললে, $\lambd[x][(x+3)]$ হলো তার আর্গুমেন্টের সঙ্গে তিন যোগ করা অপেক্ষকটির
একটি \emph{নাম}। এই রাশিতে $x$ একটি সহায়ক চলরাশি বা স্থানধারক: একই
অপেক্ষককে $\lambd[y][(y + 3)]$ দিয়েও সমানভাবে নির্দেশ করা যায়। অন্য
পরামিতি থাকলেও এই লিপি কাজ করে। যেমন, ধরা যাক $g(x, y)$ দুই চলরাশির একটি
অপেক্ষক এবং $k$ একটি স্বাভাবিক সংখ্যা। তাহলে $\lambd[x][g(x,k)]$ হলো সেই
অপেক্ষক, যা যেকোনো~$x$-কে~$g(x, k)$-তে পাঠায়।

সাংকেতিক রাশি থেকে অপেক্ষক সংজ্ঞায়িত করার এই পদ্ধতিকে
\emph{ল্যাম্বডা-বিমূর্তন} বলে। ল্যাম্বডা-বিমূর্তনের বিপরীত দিকটি হলো
\emph{অপেক্ষকের প্রয়োগ}: কোনো অপেক্ষক~$f$ (ধরা যাক, স্বাভাবিক সংখ্যাগুলিতে
সংজ্ঞায়িত) থাকলে তাকে যেকোনো মানে, যেমন 2-এ, \emph{প্রয়োগ} করা যায়।
প্রচলিত লিপিতে অবশ্য ফলটিকে~$f(2)$ লিখি।

ল্যাম্বডা-বিমূর্তন ও প্রয়োগকে একত্র করলে কী ঘটে? তখন বিমূর্ত করা চলরাশির
স্থানে প্রয়োগিত রাশিটি ``বসিয়ে'' পাওয়া রাশিকে সরল করা যায়। যেমন,
\[
(\lambd[x][(x + 3)])(2)
\]
কে~$2 + 3$-এ সরল করা যায়।

এ পর্যন্ত আমরা প্রচলিত ধারণাগুলির জন্য নতুন লিপি প্রবর্তন ছাড়া আর কিছুই
করিনি। কিন্তু ল্যাম্বডা ক্যালকুলাস সেট-তাত্ত্বিক দৃষ্টিভঙ্গি থেকে আরও মৌলিক
বিচ্যুতি ঘটায়। এই কাঠামোতে:
\begin{enumerate}
\item সবকিছুই একটি অপেক্ষক নির্দেশ করে।
\item ল্যাম্বডা-বিমূর্তন ব্যবহার করে অপেক্ষক সংজ্ঞায়িত করা যায়।
\item যেকোনো কিছুকে অন্য যেকোনো কিছুর উপর প্রয়োগ করা যায়।
\end{enumerate}
যেমন, $F$ ল্যাম্বডা ক্যালকুলাসের একটি পদ হলে $F(F)$-কে সর্বদা অর্থপূর্ণ
ধরা হয়। এই উদার কাঠামোকে \emph{টাইপবিহীন} ল্যাম্বডা ক্যালকুলাস বলে;
এখানে ``টাইপবিহীন'' মানে ``কোন জিনিসকে কোন জিনিসের উপর প্রয়োগ করা যাবে,
তার কোনো বাধা নেই।''

\begin{digress}
\emph{টাইপযুক্ত} ল্যাম্বডা ক্যালকুলাসও আছে; এটি টাইপবিহীন সংস্করণটির
একটি গুরুত্বপূর্ণ প্রকরণ। টাইপযুক্ত ল্যাম্বডা ক্যালকুলাস অনেক দিক থেকে
টাইপবিহীনটির অনুরূপ হলেও তাকে চিরায়ত সেট-তাত্ত্বিক কাঠামোর সঙ্গে সামঞ্জস্যপূর্ণ
করা অনেক সহজ, এবং তার কয়েকটি ধর্ম একেবারেই আলাদা।

তাত্ত্বিক কম্পিউটার বিজ্ঞান এবং প্রোগ্রামিং ভাষার নকশায় ল্যাম্বডা ক্যালকুলাসের
গবেষণা কেন্দ্রীয় গুরুত্ব পেয়েছে। ১৯৫০-এর দশকে জন ম্যাকার্থির নকশা করা
লিস্প এই ধারণাগুলি দ্বারা প্রভাবিত ভাষার একটি গোড়ার উদাহরণ।
\end{digress}

\end{document}
